% Part: sets-functions-relations
% Chapter: relations
% Section: special-properties

\documentclass[../../../include/open-logic-section]{subfiles}

\begin{document}

\olfileid[pt]{sfr}{rel}{prp}
\olsection{Propriedades Especiais das Relações}

\begin{intro}
Alguns tipos de relações revelam-se tão comuns que receberam nomes especiais.
Por exemplo, $\le$ e~$\subseteq$ relacionam
os seus respectivos domínios (digamos, $\Nat$ no caso de~$\le$ e
$\Pow{A}$ no caso de~$\subseteq$) de modos semelhantes. Para precisar
exatamente como estas relações são semelhantes, e como diferem,
as categorizamos segundo algumas propriedades especiais que as relações
podem ter. Verifica-se que (combinações de) algumas destas propriedades
especiais são especialmente importantes: ordens e relações de equivalência.
\end{intro}

\begin{defn}[Reflexividade]
Uma relação $R \subseteq A^2$ é \emph{reflexiva} se e só se, para todo $x \in
A$, $Rxx$.
\end{defn}

\begin{defn}[Transitividade]
Uma relação $R \subseteq A^2$ é \emph{transitiva} se e só se, sempre que $Rxy$
e $Ryz$, então também $Rxz$.
\end{defn}

\begin{defn}[Simetria]
Uma relação~$R \subseteq A^2$ é \emph{simétrica} se e só se, sempre que
$Rxy$, então também~$Ryx$.
\end{defn}

\begin{defn}[Antissimetria]
Uma relação~$R \subseteq A^2$ é \emph{antissimétrica} se e só se, sempre que tanto
$Rxy$ quanto $Ryx$, então $x=y$ (ou, noutras palavras: se $x\neq y$ então
ou $\lnot Rxy$ ou $\lnot Ryx$).
\end{defn}

\begin{explain}
Numa relação simétrica, $Rxy$ e $Ryx$ verificam-se sempre em conjunto, ou
nenhuma se verifica. Numa relação antissimétrica, a única forma de $Rxy$
e $Ryx$ se verificarem em conjunto é que $x = y$. Note-se que isto não
\emph{exige} que $Rxy$ e $Ryx$ se verifiquem quando $x = y$, apenas que isso
não está excluído. Assim, uma relação antissimétrica pode ser reflexiva, mas
não é o caso que toda relação antissimétrica seja
reflexiva. Note-se também que ser antissimétrica e ser meramente não simétrica
são condições diferentes. De fato, uma relação pode ser ao mesmo tempo
simétrica e antissimétrica (por exemplo, a relação identidade o é).
\end{explain}

\begin{defn}[Conexidade]
Uma relação $R \subseteq A^2$ é \emph{conexa} se para todos $x,y\in
A$, se $x \neq y$, então ou $Rxy$ ou~$Ryx$.
\end{defn}

\begin{prob}
Dê exemplos de relações que sejam (a) reflexivas e simétricas mas
não transitivas, (b) reflexivas e antissimétricas, (c) antissimétricas,
transitivas, mas não reflexivas, e (d) reflexivas, simétricas e
transitivas. Não use relações sobre números ou conjuntos.
\end{prob} 
 
\begin{defn}[Irreflexividade]
Uma relação $R \subseteq A^2$ diz-se \emph{irreflexiva} se, para todo $x \in
A$, não $Rxx$. 
\end{defn}

\begin{defn}[Assimetria]
Uma relação $R \subseteq A^2$ diz-se \emph{assimétrica} se para nenhum par $x,y\in
A$ se tem tanto $Rxy$ quanto~$Ryx$. 
\end{defn}

Note-se que se $A \neq \emptyset$, então nenhuma relação irreflexiva em~$A$
é reflexiva e toda relação assimétrica em~$A$ é também
antissimétrica. Contudo, há $R \subseteq A^2$ que não são
reflexivas nem irreflexivas, e há relações antissimétricas
que não são assimétricas. 

\end{document}
